%% =============================================================================
%%  CHAPTER 3 — Methodology
%%  Reusable template. Edit the content below for your topic.
%%  Variables are loaded from variables.tex in the same directory.
%% =============================================================================

\chapter{Methodology}

\section{Introduction}

%% ── Overview ─────────────────────────────────────────────────────────────────
%% Introduce the methodology chapter.
%% Describe the overall approach and workflow.

% This chapter describes the methodology used to [achieve objective].
% The approach consists of [N] main stages: [stage 1], [stage 2], [stage 3].
% Figure~\ref{fig:workflow} provides an overview of the pipeline.

%% ── Workflow Figure ──────────────────────────────────────────────────────────
% \begin{figure}[htb!]
% \centering
% \includegraphics[width=\textwidth]{workflow.png}
% \caption{Overview of the methodological workflow}
% \label{fig:workflow}
% \end{figure}

\section{Data Collection}

%% ── Data Sources ─────────────────────────────────────────────────────────────
%% Describe where the data comes from, how it was collected, and any filters.

% \subsection{Primary Data Sources}
% \subsection{Secondary Data Sources}
% \subsection{Data Filtering Criteria}

\section{Data Preprocessing}

%% ── Preprocessing Steps ──────────────────────────────────────────────────────
%% Describe cleaning, normalization, transformation steps.

% \subsection{Cleaning}
% \subsection{Normalization}
% \subsection{Feature Engineering}

\section{[Core Method / Approach]}

%% ── Main Method ──────────────────────────────────────────────────────────────
%% Describe the core methodology in detail.

% \subsection{[Component 1]}
% \subsection{[Component 2]}
% \subsection{[Component 3]}

\section{Experimental Setup}

%% ── Environment & Configuration ──────────────────────────────────────────────
%% Describe hardware, software, parameters.

% \subsection{Hardware / Software Environment}
% \begin{itemize}
%     \item CPU/GPU specifications
%     \item Programming language and version
%     \item Libraries and dependencies
% \end{itemize}

% \subsection{Hyperparameters}
% \subsection{Evaluation Metrics}

% \def\metricExample{0.95}  % Example: target accuracy

\section{Validation Strategy}

%% ── Validation ───────────────────────────────────────────────────────────────
%% How do you validate your approach? Cross-validation? Ablation? Human eval?

% \subsection{Cross-Validation}
% \subsection{Ablation Studies}
% \subsection{Human-in-the-Loop Verification}

\section{Summary}

%% ── Chapter Summary ──────────────────────────────────────────────────────────
% This chapter described the methodology comprising [N] stages:
% [data collection], [preprocessing], [core method], [validation].
% Chapter~4 presents the results of applying this methodology.
